\documentclass[../DoAn.tex]{subfiles}
\begin{document}

Đồ án làm ra một hệ thống chấm công nhận diện khuôn mặt chạy trên ESP32-S3, gồm máy chấm công, máy chủ quản trị nhân sự
trên VPS, dashboard web và cổng nhân viên trên điện thoại. Máy chấm công phát hiện, chống giả mạo và nhận diện khuôn mặt
ngay trên chip, ảnh khuôn mặt không rời máy, và máy vẫn chấm công khi mất mạng. Máy chủ đi từ lượt chấm tới ngày công,
đơn từ, kỳ lương và phiếu lương có thuế thu nhập cá nhân và bảo hiểm theo các mức pháp lý năm 2026.

Phần trên chip được đo trên board thật. Một lượt ba mô hình mất 353,6~ms lúc rảnh và 406,2~ms khi có xem trước, nhanh hơn
TFLite Micro dùng esp-nn 10,51 lần. Khoảng cách này đến từ luật lượng tử của ESP-DL và các phép toán đồ án tự viết; riêng
các bản vá giữ nguyên số nguyên đã giảm 25,0\% thời gian mà đầu ra không đổi một bit. Đáy RAM nội là 28.675~B, trên mốc
24~KB. Bộ phát hiện đạt AP 0,9544 trên mặt từ 38~px. Nhánh chống giả mạo chặn 25/25 khung giả và giữ 59/62 mặt thật trên
87 khung OV5640, vì mô hình được chọn bằng chính đường ảnh của chip: bản tự huấn luyện chỉ đạt AUC 0,6473 trên
85 khung cùng loại.

Đường lượt chấm ghi flash trước, gửi QoS 1, chống trùng bằng ràng buộc unique trên cặp máy và mã lượt, và xả 604 bản ghi
tồn qua TLS không lỗi. Firmware và gói mô hình đều có hai khe, bản mới tự xác nhận sau 30~s chạy thử, và chín bản kê
khai sai đều bị từ chối đúng tên. Backend qua 971/971 ca e2e, gồm phân quyền theo dòng dữ liệu, tính công và tính
lương.

Hệ thống còn năm hạn chế. (i)~NF1 chỉ đạt lúc rảnh: ở điểm vận hành, bus PSRAM dùng chung với đường xem trước kéo một lượt
lên 406,2~ms. (ii)~Nhận diện sụt 3,7\% TAR trên CFP-FP sau lượng tử hóa dù đã qua QAT, nên máy từ chối oan nhiều hơn với
mặt nghiêng; trọng số nhận diện học trên MS-Celeb-1M chỉ được dùng cho nghiên cứu. (iii)~Tập OV5640 đã dùng để chọn mô
hình, cả 25 khung giả đều là màn hình điện thoại, chưa có khung ảnh in có nhãn, và FAR, FRR với nhiều người chưa đo. (iv)~Trên
VPS 1 vCPU, 2~GB với dữ liệu 5.000 nhân viên, hệ thống chịu được 10 phiên nhưng sập ở 20 phiên, và lượt lương 4.776 người
thất bại vì BullMQ mất khóa; ba nguyên nhân ở Mục~\ref{section:5.6.2} chưa được sửa. (v)~Máy còn lắp trên breadboard,
mạch đế chưa in, chưa có vỏ và thanh chắn; khi Flash Encryption và Secure Boot chưa bật, người có máy trong tay đọc được JWT.

Việc trước mắt là sửa ba lỗi của máy chủ rồi đo lại trên VPS: thêm chỉ mục \texttt{WHERE review IS NOT NULL}, với câu đếm
lượt chờ duyệt còn 0,08~ms trên dữ liệu đo; thay healthcheck của EMQX bằng một phép kiểm không dựng máy ảo Erlang; sửa việc
nền tính lương để BullMQ không mất khóa khi lượt chạy kéo dài. Tiếp theo là thu khung OV5640 có nhãn của ảnh in và thẻ,
đo FAR, FRR với nhiều người trên một tập không dùng để chọn mô hình, và chạy các phép thử còn lại ở
Bảng~\ref{tab:chua_do}. Nhánh nhận diện cần huấn luyện lại trên dữ liệu được phép dùng thương mại, và thử QAT dài hơn
2.000 bước để đưa TAR lên cửa 0,822.

Về phần cứng, mạch đế cần được in, máy cần vỏ và thanh chắn thật, rồi đo dòng tiêu thụ và thời gian chạy bằng pin. Bước
cuối trước khi lắp đặt là bật Flash Encryption, Secure Boot và mã hóa NVS. Về nghiệp vụ, ca làm qua nửa đêm chưa có thiết
kế và hiện bị từ chối với mã \texttt{SHIFT\_CROSSES\_MIDNIGHT}, nên doanh nghiệp có ca vắt qua 0 giờ chưa dùng được hệ thống.

\end{document}
